We measured adaptation and retention during continued SAC training under sensor faults, with eight seeds per controller, two tasks, a clean-continuation control and pre-specified paired tests. In HalfCheetah, delay separated the two outcomes: the reactive controller adapted while losing clean performance, and the transformer lost clean performance without adapting. Retaining half a batch of earlier clean experience reduced every clean-condition loss, at the cost of part of the reactive controller's delay gain, and the critic diagnostic is consistent with values blended across the two sensing channels. In Ant, standing sources tolerated both faults, so that task probed continued training more than adaptation. Next steps are matched ablations that separate memory from capacity, explicit retention objectives, and source recipes that produce locomotion.
